\chapter{Symmetry in Quantum Physics}

\section{Symmetry Groups in Quantum Physics}

\textbf{Symmetry transformations} are operations carried out on the devices which prepare the states and the apparatuses which measure the observables of a quantum system having the property of not changing the results of experiments. For instance, space translations are symmetry transformations since the uniform translation of all preparations devices and measuring apparatuses does not change the results of experiments due to the homogeneity of space. Likewise, space rotations are symmetry transformations since the uniform rotation of all preparations devices and measuring apparatuses does not change the results of experiments due to the isotropy of space. In relativity, Lorentz boosts are symmetry transformations since as the result of experiments are the same in a given inertial frame and any inertial frame moving at uniform speed with respect to the first because the isotropy of spacetime.

In general, symmetry transformations gather together in families with special properties. Let \(G\) be one such family. Carrying out successively two symmetry transformations \(g\), \(g'\) of \(G\) yields a third symmetry transformation \(g' \cdot g\) of \(G\) called the \textbf{product} of \(g\), \(g'\). The effect of a symmetry transformation \(g\) of \(G\) can be undone by another symmetry transformation \(g^{-1}\) of \(G\) called the \textbf{inverse} of \(g\). There exists a distinguished symmetry transformation \(e\) of \(G\) that has no effect called \textbf{neutral}. The multiplication and inversion operations and neutral element of \(G\) have the following properties:
\begin{itemize}
    \item \textbf{Associativity}: for any three symmetry transformations \(g\), \(g'\), \(g''\) of \(G\) we have
          \begin{equation}
              (g'' \cdot g') \cdot g = g'' \cdot (g' \cdot g).
              \label{eq:group_associativity_definition}
          \end{equation}
    \item \textbf{Neutrality}: there exists a symmetry transformation \(e\) of \(G\) such that for any symmetry transformation \(g\) of \(G\) we have
          \begin{equation}
              e \cdot g = g \cdot e = g.
              \label{eq:group_neutrality_definition}
          \end{equation}
    \item \textbf{Invertibility}: for any symmetry transformation \(g\) of \(G\) there exists a symmetry transformation \(g^{-1}\) of \(G\) such that
          \begin{equation}
              g^{-1} \cdot g = g \cdot g^{-1} = e.
              \label{eq:group_invertibility_definition}
          \end{equation}
\end{itemize}
These relations characterize \(G\) as an algebraic structure called a \textbf{group}. Since \(G\) is constituted by symmetry transformations, \(G\) is called more specifically a \textbf{symmetry group}. A group \(G\) is called \textbf{finite} if it is constituted by a finite number of elements, else it is called \textbf{infinite}. A group \(G\) is called \textbf{Abelian} if multiplication turns out to be commutative, that is
\begin{equation}
    g \cdot g' = g' \cdot g, \quad \forall g, g' \in G.
    \label{eq:group_abelian_definition}
\end{equation}
Else, it is called non Abelian. The groups encountered in quantum physics are bot finite and infinite and Abelian and non Abelian.

\begin{example}[Space Translation Symmetry]
    [...]\TODO{Complete the example}
\end{example}

\begin{example}[Space Rotation Symmetry]
    [...]\TODO{Complete the example}
\end{example}

\subsection{Symmetry Groups and Their Operator Realization}

In quantum theory, states and observables of a system are represented respectively by normalized vectors defined up to a multiplicative phase in a Hilbert space \(\mathcal{H}\) and selfadjoint operators in the operator algebra \(L(\mathcal{H})\). The invariance of experimental results under the action of a symmetry group reduces to the invariance of expectation values: if \(\Psi\) is a state, \(O\) is an observable and \(g\) is a symmetry transformation of \(G\), \(^g \Psi\) is the transformed state of \(\Psi\) and \(^g O\) is the transformed observable of \(O\), one must have
\begin{equation}
    \bra{^g \Psi} ^g O \ket{^g \Psi} = \bra{\Psi} O \ket{\Psi}.
    \label{eq:group_invariance_expectation_value_definition}
\end{equation}

Let us now analyze the implications of the invariance of expectation values under the action of a symmetry group \(G\) with a simple example.

\begin{example}[``yes–no'' observables]
    As is well–known, there is a class of elementary observables, the “yes–no” observables, defined by taking the value 1, 0 according to whether the system is in a certain state \(\Phi\) or not. The associated operator is given by
    \[
        O_{\Phi} \Psi = \Phi \bra{\Phi} \ket{\Psi}.
    \]
    It is natural to assume that under the symmetry transformation \(g\) of \(G\) the observable \(O_{\Phi}\) transforms accordingly with the state \(\Phi\) as follows:
    \[
        ^g O_{\Phi} = O_{^g \Phi}.
    \]
    The expectation value of \(O_{\Phi}\) in the state \(\Psi^{\prime}\) is given by the transition probability from \(\Psi^{\prime}\) to \(\Phi\):
    \[
        \bra{\Psi^{\prime}} O_{\Phi} \ket{\Psi^{\prime}} = |\bra{\Phi} \ket{\Psi^{\prime}}|^2,
    \]
    such that the \textbf{invariance of expectation values} under the action of the symmetry group \(G\) implies that
    \[
        |\bra{^g \Phi} \ket{^g \Psi^{\prime}}|^2 = |\bra{\Phi} \ket{\Psi^{\prime}}|^2.
    \]
    This means that the symmetry transformation \(g\) of \(G\) must be implemented by a unitary or antiunitary operator \(U_g\) on the Hilbert space \(\mathcal{H}\) of the system such that
    \[
        \ket{^g \Psi} = U_g \ket{\Psi}, \quad ^g O = U_g O U_g^{-1}.
    \]
\end{example}

In this example we have highlighted the importance of the \textbf{Wigner's theorem}:
\begin{theorem}[Wigner's theorem]
    Let \(\mathcal{H}\) be a Hilbert space and let \(G\) be a symmetry group of a quantum system represented on \(\mathcal{H}\). Then, for each symmetry transformation \(g\) of \(G\), there exists a unitary or antiunitary operator \(U_g\) on \(\mathcal{H}\) such that
    \begin{equation}
        \begin{dcases}
            ^g \Psi = U(g) \Psi,        \\
            ^g O = U_g O U_g^{\dagger}, \\
        \end{dcases}
        \label{eq:unitary_representation_symmetry_group_definition}
    \end{equation}
    since \(U_g\) is unitary or antiunitary, \(U_g^{-1} = U_g^{\dagger}\), and so we can write the transformation of observables as \(O \mapsto U_g O U_g^{\dagger}\). This implies that the expectation value of any observable \(O\) in any state \(\Psi\) is invariant under the action of the symmetry group \(G\).
\end{theorem}
This is a fundamental result in quantum theory since it allows us to represent symmetry transformations as unitary or antiunitary operators on the Hilbert space of the system.

Recall that a unitary operator \(U\) on a Hilbert space \(\mathcal{H}\) is a linear operator that preserves the Hilbert inner product,
\[
    \begin{aligned}
        U (\Phi\,a ,\, \Phi'\,a')  & = U (\Phi) a,\, U(\Phi')a^{\prime} \\
        \bra{U \Phi} \ket{U \Phi'} & = \bra{\Phi} \ket{\Phi'},
    \end{aligned}
\]
while an antiunitary operator \(A\) on a Hilbert space \(\mathcal{H}\) is an antilinear operator that preserves the Hilbert inner product,
\[
    \begin{aligned}
        A (\Phi\,a ,\, \Phi'\,a')  & = A (\Phi) a^*,\, A(\Phi')a^{\prime *} \\
        \bra{A \Phi} \ket{A \Phi'} & = \bra{\Phi} \ket{\Phi'}^*.
    \end{aligned}
\]
in both cases the adjoint operator \(U^{\dagger}\) (or \(A^{\dagger}\)) of a unitary (or antiunitary) operator \(U\) (or \(A\)) are equal to the inverse operator \(U^{-1}\) (or \(A^{-1}\)):
\[
    \begin{dcases}
        \bra{\Psi^{\prime}}\ket{U \Psi} = \bra{U^{\dagger} \Psi^{\prime}} \ket{\Psi}, \\
        \bra{\Psi^{\prime}}\ket{A \Psi} = (\bra{A^{\dagger} \Psi^{\prime}} \ket{\Psi})^*,
    \end{dcases} \quad \implies \quad \begin{dcases}
        U^{-1} = U^{\dagger}, \\
        A^{-1} = A^{\dagger},
    \end{dcases}
\]
respectively for the linear and antilinear case. In what follows, we leave aside the antiunitary case, which occurs only for symmetry transformations involving time reversal, and concentrate on the unitary one. Letting the symmetry transformation \(g\) vary in the group \(G\), we obtain a family of unitary operators \(U(g)\) contained in the operator algebra \(L(\mathcal{H})\).

Let us come back to relation \eqref{eq:unitary_representation_symmetry_group_definition}\(_1\). We recall again that state vectors are normalized vectors defined up to a multiplicative phase factor. The content of \eqref{eq:unitary_representation_symmetry_group_definition}\(_1\) is that phase choice can be adjusted so that for any symmetry transformation \(g\) the state correspondence \(\Psi \mapsto ^g \Psi\) is codified by a unitary operator \(U(g)\). This however does not completely fix \(U(g)\): we are still free to replace \(U(g)\) by
\begin{equation}
    \hat{U}(g) = \alpha(g) U(g),
    \label{eq:unitary_representation_phase_redefinition}
\end{equation}
where \(\alpha(g)\) is a complex number of unit modulus, which is called a \textbf{phase factor}. This involves a phase redefinition
\[
    ^g \Psi \to \alpha(g) \, ^g \Psi,
\]
in \eqref{eq:unitary_representation_symmetry_group_definition}\(_1\), which does not change the physical content of the transformation since states are defined up to a multiplicative phase factor. By a suitable choice of the phase factors \(\alpha(g)\) many relations simplyfy.

It is natural and possible to require
\[
    ^e \Psi = \Psi, \quad \forall \Psi \in \mathcal{H},
\]
so that the neutral element \(e\) of the symmetry group \(G\) is represented by the identity operator
\[
    U(e) = \mathbb{I},
\]
on the Hilbert space \(\mathcal{H}\) of the system. With this choice, we can now analyze the implications of the group properties of \(G\) on the family of unitary operators \(U(g)\). Considering two symmetry transformations \(g\), \(g'\) of \(G\), and a physical state \(\Psi\), we have \(^{g^{\prime}}(^g \Psi)\) and \(^{g^{\prime} \cdot g} \Psi\) defined only up to a conventional choice of a multiplicative \textbf{phase factor} as seen above, one cannot expect that a relation of the form \(^{g^{\prime}}(^g \Psi) = ^{g^{\prime} \cdot g} \Psi\) to hold strictly in general, but only up to some phase factor
\[
    \gamma(g, g^{\prime}, \Psi) \in \mathbb{C}, \quad |\gamma(g, g^{\prime}, \Psi)| = 1,
\]
depending on the symmetry transformations \(g\), \(g'\) and the state \(\Psi\). Thus we can write
\[
    ^{g^{\prime}}(^g \Psi) = ^{g^{\prime} \cdot g} \Psi \, \gamma(g, g^{\prime}, \Psi).
\]
It can be proven that the phase factor \(\gamma(g, g^{\prime}, \Psi)\) can be chosen to be independent of the state \(\Psi\), thus obtaining a relation of the form
\begin{equation}
    U(g \cdot g^{\prime}) = \gamma(g, g^{\prime}) U(g) U(g^{\prime}),
    \label{eq:projective_representation_symmetry_group_definition}
\end{equation}
which is called a \textbf{projective representation} of the symmetry group \(G\). If the phase factor \(\gamma(g, g^{\prime})\) can be chosen to be 1 for all \(g\), \(g'\) in \(G\), then we have a \textbf{unitary representation} of the symmetry group \(G\).

The phase factors \(\gamma(g, g^{\prime})\) appearing in the projective representation of the symmetry group \(G\) are not arbitrary, but must satisfy some consistency conditions, for example the \textit{2-cocyclicity}: given three symmetry transformations \(g\), \(g'\), \(g''\) of \(G\), the following relation must hold:
\begin{equation}
    \gamma(g, g' \cdot g'') \gamma(g', g'') = \gamma(g \cdot g', g'') \gamma(g, g'),
    \label{eq:2-cocyclicity_definition}
\end{equation}
where functions \(\gamma : G \times G \to U(1)\) satisfying the 2-cocyclicity condition are called \textbf{2-cocycles} of \(G\). The 2-cocycle condition is a consequence of the associativity property of the group multiplication of \(G\).

The \(G\) 2-cocycle \(\gamma(g, g')\) enjoys a few special properties, which albeit technical are often useful at various points:
\begin{equation}
    \begin{aligned}
        \gamma(e, g)      & = \gamma(g, e) = 1,  \\
        \gamma(g, g^{-1}) & = \gamma(g^{-1}, g),
    \end{aligned}
    \label{eq:2-cocycles_properties}
\end{equation}
for any symmetry transformation \(g\) of \(G\). The first property is a consequence of the neutrality property of the group multiplication of \(G\), while the second one is a consequence of the invertibility property of the group multiplication of \(G\).

Relation \eqref{eq:2-cocyclicity_definition} entails also that, for any symmetry transformation \(g\), the operators \(U(g^{-1})\) and \(U(g)^{-1}\) differ by a \(g\)-dependent phase factor. Specifically, the relation between them takes the form
\[
    U(g^{-1}) = \varpi(g) U(g)^{-1}, \quad \varpi(g) = \gamma(g, g^{-1})^{-1},
\]

In general it is not possible to absorb the phase \(\gamma(g^{\prime},g)\) by redefining the operators \(U(g)\) into the operators \(\hat{U}(g)\) given in \eqref{eq:unitary_representation_phase_redefinition} with \(\alpha(g)\) a suitable phase function. By doing so \(\gamma(g^{\prime},g)\) gets replaced by
\begin{equation}
    \hat{\gamma}(g, g') = \gamma(g, g') \alpha(g \cdot g') \alpha(g)^{-1} \alpha(g')^{-1},
    \label{eq:2-cocycle_phase_factor_definition}
\end{equation}
and in general this cannot be made equal to 1 by a judicious choice of \(\alpha(g)\). A phase function \(\gamma(g^{\prime},g)\) of the form
\begin{equation}
    \gamma(g, g') = \alpha(g \cdot g') \alpha(g)^{-1} \alpha(g')^{-1},
    \label{eq:2-coboundary_definition}
\end{equation}
for some function \(\alpha : G \to U(1)\) is called a \textbf{2-coboundary} of \(G\) (which is a trivial 2-cocycle).

The G 2–cocycle \(\hat{\gamma}(g^{\prime},g)\) in \eqref{eq:2-cocycle_phase_factor_definition} can be rendered equal to 1 precisely when the G 2–cocycle \(\gamma(g^{\prime},g)\) is a G 2–coboundary of the form \eqref{eq:2-coboundary_definition}. This may or may not happen. The study of these matters is the object of a branch of algebra known as group cohomology. Some special results will be obtained below. At any rate, if \(\hat{\gamma}(g^{\prime},g)\) can be made equal to 1, the representation \(\hat{U}(g)\) becomes
genuinely unitary, so that
\[
    \hat{U} (g \cdot g') = \hat{U}(g) \hat{U}(g'),
\]
for all symmetry transformations \(g\), \(g'\) of \(G\). Then the unitary family \(U(g)\) forms a genuine \textbf{unitary representation} of the symmetry group \(G\).


\section{Dynamical Symmetry Groups}

\subsection{Energy Eigenvalues Classification}



\section{Lie Groups and Lie Algebras}

A Lie group \(G\) is an infinite group that can be parametrized in a neighborhood of the neutral element by continuous parameters \(\theta^a\), \(a = 1, \ldots, d\), organized as a vector \(\bs{\theta}\), in one–to–one fashion:
\[
    \begin{dcases}
        \bs{\theta} \in D \subseteq \mathbb{R}^d, \\
        g \, : \, D \mapsto G,                    \\
        \bs{\theta} \mapsto g(\bs{\theta}),
    \end{dcases}
\]
where \(D\) is an open neighborhood of the origin in \(\mathbb{R}^d\) (\(\mathbf{0}\in D\) and \(g(\mathbf{0}) \in G\)). The number \(d\) of parameters is called the dimension of \(G\), seen as a differential map. The group operations of multiplication and inversion are smooth functions of the parameters.

If the parametrization is smooth, as we assume, \(G\) can be thought geometrically as a group manifold. Below, we treat \(G\) as a matrix group for definiteness.

We assume that 0 belongs to the parameter space and that
\[
    g(\mathbf{0}) = 1,
\]
which can be read as a normalizing condition. By virtue of the parametrization, for any pair of parameter vectors \(\bs{\theta}\), \(\bs{\theta}'\), there exists a third parameter vector \(\bs{\theta}''\) such that
\[
    g(\bs{\theta}') g(\bs{\theta}) = g(\bs{\theta}'').
\]
As the parametrization is one–to–one,\footnote{This means that each element of the group corresponds to a unique set of parameters: \(g(\bs{\theta})=g(\bs{\theta}') \implies \bs{\theta} = \bs{\theta}'\).} \(\bs{\theta}''\) is also uniquely determined by \(\bs{\theta}\), \(\bs{\theta}'\).
So, it can be expressed as a function of them
\[
    \bs{\theta}'' = m(\bs{\theta}', \bs{\theta}),
\]
which is a function of the \textbf{group multiplication}. The function \(m\) is smooth as the group multiplication is smooth. The inverse of \(g(\bs{\theta})\) is given by
\[
    g(\bs{\theta})^{-1} = g(\bs{\theta}'),
\]
for some \(\bs{\theta}' = j(\bs{\theta})\), which is unique since the parameterization is one-to-one. The function \(j\), implementing the group inversion, is smooth as the group inversion is smooth.

We have
\[
    \begin{dcases}
        m : D \times D \to D, \\
        j : D \to D.
    \end{dcases}
\]

[...]

These functions have to respect some properties, which are the group multiplication properties: associativity, invertibility and neutrality. Thus given three parameter vectors \(\bs{\theta}\), \(\bs{\theta}'\), \(\bs{\theta}''\), the following relations must hold:
\[
    \begin{aligned}
        \left( g(\bs{\theta}) g(\bs{\theta}') \right) g(\bs{\theta}'') & = g(\bs{\theta}) \left( g(\bs{\theta}') g(\bs{\theta}'') \right), \\
        g(m(\bs{\theta}, \bs{\theta}')) g(\bs{\theta}'')               & = g(\bs{\theta})g(m(\bs{\theta}', \bs{\theta}'')),                \\
        g(m(m(\bs{\theta}, \bs{\theta}'), \bs{\theta}''))              & = g(m(\bs{\theta}, m(\bs{\theta}', \bs{\theta}''))),
    \end{aligned}
\]
since the parametrization is one–to–one, these relations
    [\dots ]
and so we find the three properties of the group multiplication functions:
\begin{equation}
    \begin{aligned}
        m(m(\bs{\theta}, \bs{\theta}'), \bs{\theta}'') & = m(\bs{\theta}, m(\bs{\theta}', \bs{\theta}'')), \\
        m(\bs{\theta}, j(\bs{\theta}))                 & = m(j(\bs{\theta}), \bs{\theta}) = \mathbf{0},    \\
        m(\bs{\theta}, \mathbf{0})                     & = m(\mathbf{0}, \bs{\theta}) = \bs{\theta}.
    \end{aligned}
\end{equation}

The idea now is to try and understand the structure of \(G\) by looking at the behavior of \(g(\bs{\theta})\) in a neighborhood of the neutral element. This is done by Taylor expanding it in power series around \(\mathbf{0}\):
\[
    g(\bs{\theta}) = g(\mathbf{0}) + \frac{\partial g}{\partial \theta^a} \bigg|_{\mathbf{0}} \theta^a + \frac{1}{2} \frac{\partial^2 g}{\partial \theta^a \partial \theta^b} \bigg|_{\mathbf{0}} \theta^a \theta^b + \mathrm{O}(\theta^3).
\]
Now computing individual terms of the expansion, we have
\[
    g(\bs{0}) = 1, \quad t_a = \frac{\partial g}{\partial \theta^a} \bigg|_{\mathbf{0}}, \quad t_{ab} = \frac{\partial^2 g}{\partial \theta^a \partial \theta^b} \bigg|_{\mathbf{0}},
\]
so that
\[
    g(\bs{\theta}) = 1 + t_a \theta^a + \frac{1}{2} t_{ab} \theta^a \theta^b + \mathrm{O}(\theta^3).
\]
Let us highlight the fact that the coefficients \(t_a\), \(t_{ab}\) are matrices, as \(g(\bs{\theta})\) is a matrix group. \(t_{ab}\) in particular is symmetric in its indices \(a\), \(b\) since the order of differentiation does not matter.

We can now think of \(G\) as an \textbf{hypersurface in the space of matrices}, containing the identity matrix. In the neighborhood of the identity, there is \(D\), subset of our hypersurface, which is diffeomorphic to an open neighborhood of the origin in \(\mathbb{R}^d\).

If we vary \(g(0,\dots,0,\theta^a,0,\dots,0)\), we get a curve in \(D\) passing through the identity (when \(\theta^a = 0\)). The tangent vector to this curve at the identity is given by \(t_a = \frac{\partial g}{\partial \theta^a} \big|_{\mathbf{0}}\). Varying all parameters \(\theta^a\) we get \(d\) curves in \(D\) passing through the identity and their tangent vectors at the identity are given by \(t_a\), \(a = 1, \ldots, d\). These tangent vectors are \textbf{linearly independent} since the parametrization is one–to–one. Thus, they span a vector space of dimension \(d\).

Thus the tangent space to \(G\) at the identity is more than a vector space of dimension \(d\) (spanned by the matrices \(t_a\), \(a = 1, \ldots, d\), costituing a basis). We denote this tangent space by \(\mathfrak{g}\) and call it the Lie algebra of \(G\).
It is a vector space of dimension \(d\) and it is closed under the commutator of matrices. The Lie algebra \(\mathfrak{g}\) encodes all the information about the local structure of \(G\) around the identity element.

We could now think to Taylor expand the group multiplication functions \(m\), \(j\) around \(\mathbf{0}\) and try to find the structure of \(G\) by looking at the properties of the coefficients of the expansion. Starting from the associativity property, we find that the coefficients of the expansion of \(m\) must satisfy
\[
    m^a(\bs{\theta}, \bs{\theta}') = \theta^a + \theta'^a + m^a_{bc} \theta^b \theta'^c + \mathrm{O}(\bs{\theta}'^2, \bs{\theta}) + \mathrm{O}(\bs{\theta}^2, \bs{\theta}'),
\]
while for the inversion property we find
\[
    j^a(\bs{\theta}) = -\theta^a + j^a_{bc} \theta^b \theta^c + \mathrm{O}(\bs{\theta}^3).
\]
If we now recall the Taylor expansion of \(g(\bs{\theta})\) around \(\mathbf{0}\)
\[
    g(\mathbf{0}) = 1 + t_a \theta^a + \frac{1}{2} t_{ab} \theta^a \theta^b + \mathrm{O}(\theta^3),
\]
we can note how the coefficients \(t_a\) define another member of the Lie algebra \(\mathfrak{g}\) by defining the \textbf{commutator}, or \textbf{Lie bracket}, of two basis elements \(t_a\), \(t_b\) of \(\mathfrak{g}\) as follows:
\[
    [t_a, t_b] = t_a t_b - t_b t_a = f_{ab}^{\ \ c} t_c \in \mathfrak{g},
\]
where \(f_{ab}^{\ \ c}\) are the \textbf{structure constants} of the Lie algebra \(\mathfrak{g}\). The structure constants are \textit{antisymmetric in their lower indices} \(a\), \(b\) since the commutator is antisymmetric. These structure constants are related to the coefficients \(m^a_{\ bc}\) of the expansion of the group multiplication function \(m\) by
\[
    f^a_{\ bc} = m^a_{\ bc} - m^a_{\ cb}.
\]
If we change parametrization, the structure constants change accordingly, but the Lie algebra \(\mathfrak{g}\) remains the same.

Taking again the definition of the multiplication function
\[
    g(\bs{\theta}') g(\bs{\theta}) = g(m(\bs{\theta}', \bs{\theta})),
\]
we can Taylor expand both sides around \(\mathbf{0}\) and find the following relation between the coefficients of the expansion of \(g\) and the structure constants:
\[
    \begin{aligned}
        \left[ 1 + \theta^a t_a + \frac{1}{2} \theta^a \theta^b t_{ab} + \mathrm{O}(\theta^3) \right] \left[ 1 + \theta'^a t_a + \frac{1}{2} \theta'^a \theta'^b t_{ab} + \mathrm{O}(\theta'^3) \right] \\
        =[\dots ]
    \end{aligned}
\]
and by relating the coefficients of the expansion of both sides in \(\theta^b\), \(\theta'^a\) we find the following relations (renaming summation indices):
\[
    \begin{aligned}
        t_a t_b & = t_{ab} + m^c_{\ ab} t_c, \\
        t_b t_a & = t_{ab} + m^c_{\ ba} t_c,
    \end{aligned}
\]
and if we subtract the second from the first, we find
\[
    [t_a, t_b] = (m^c_{\ ab} - m^c_{\ ba}) t_c = f_{ab}^{\ \ c} t_c,
\]
so that the structure constants \(f_{ab}^{\ \ c}\) of the Lie algebra \(\mathfrak{g}\) are related to the coefficients \(m^c_{\ ab}\) of the expansion of the group multiplication function \(m\) by
\[
    f_{\ ab}^{c} = m^c_{\ ab} - m^c_{\ ba}.
\]

Thus to sum up:
\begin{itemize}
    \item The Lie algebra \(\mathfrak{g}\) of a Lie group \(G\) is the tangent space to \(G\) at the identity element, spanned by the matrices \(t_a\), \(a = 1, \ldots, d\).
    \item The basis elements \(t_a\) of the Lie algebra \(\mathfrak{g}\) satisfy the commutation relations
          \[
              [t_a, t_b] = f_{ab}^{\ \ c} t_c,
          \]
          where \(f_{ab}^{\ \ c}\) are the structure constants of the Lie algebra \(\mathfrak{g}\).
    \item The structure constants \(f_{ab}^{\ \ c}\) of the Lie algebra \(\mathfrak{g}\) are related to the coefficients \(m^c_{\ ab}\) of the expansion of the group multiplication function \(m\) by
          \[
              f_{\ ab}^{c} = m^c_{\ ab} - m^c_{\ ba}.
          \]
\end{itemize}

The Lie algebra, the tangent space to the Lie group at the identity, is not only a vector space, but also an algebra, since it is closed under the commutator of matrices. The Lie algebra encodes all the information about the local structure of the Lie group around the identity element. In particular, it encodes the structure constants \(f_{ab}^{\ \ c}\) which determine the commutation relations of the basis elements \(t_a\) of the Lie algebra.

The commutator of the basis elements \(t_a\) of the Lie algebra \(\mathfrak{g}\) has some properties, which are the \textbf{Jacobi identity} and the \textbf{antisymmetry} of the commutator:
\[
    \begin{aligned}
        \left[ t_a,\,t_b \right] + \left[ t_b,\,t_a \right] = 0, \\
        \left[ t_a,\,\left[ t_b,\,t_c \right] \right] + \left[ t_b,\,\left[ t_c,\,t_a \right] \right] + \left[ t_c,\,\left[ t_a,\,t_b \right] \right] = 0.
    \end{aligned}
\]
These two properties can be expressed in terms of the structure constants \(f_{ab}^{\ \ c}\) as follows:
\[
    \begin{aligned}
        f_{ab}^{\ \ c} + f_{ba}^{\ \ c}                                                               & = 0, \\
        f_{ad}^{\ \ e} f_{bc}^{\ \ d} + f_{bd}^{\ \ e} f_{ca}^{\ \ d} + f_{cd}^{\ \ e} f_{ab}^{\ \ d} & = 0,
    \end{aligned}
\]
which are the \textbf{antisymmetry} of the structure constants in their lower indices and the \textbf{Jacobi identity} for the structure constants, respectively. This means that not every set of numbers \(f_{ab}^{\ \ c}\) can be the structure constants of a Lie algebra, but only those that satisfy the antisymmetry and Jacobi identity properties.

\subsection{Exponential Parametrization of Lie Groups}

Given the parametrization of a Lie group \(G\) in a neighborhood of the identity element we used to define the Lie algebra \(\mathfrak{g}\) of \(G\)
\[
    g \, : \, D \mapsto G, \quad \bs{\theta} \in D \mapsto g(\bs{\theta}) \in G,
\]
we can now try to find a global parametrization of \(G\) in terms of the elements of \(\mathfrak{g}\). If we consider another parametrization of \(G\) in a neighborhood of the identity element defined by
\[
    \tilde{g}(\theta) = \lim_{N \to \infty} g\left(\frac{\theta}{N}\right)^N,
\]
we obtain the \textbf{exponential map}, which provides a global parametrization of the Lie group \(G\) in terms of the elements of its Lie algebra \(\mathfrak{g}\). Let us analyze its behaviour in a neighborhood of the identity element. We have
\[
    g\left(\frac{\theta}{N}\right)^N = \left[ 1 + \frac{t_a \theta^a}{N} + \frac{1}{2} \left( \frac{t_a \theta^a}{N} \right)^2 + \mathrm{O}\left(\frac{1}{N^3}\right) \right]^N \sim \left(1 +  \frac{t_a \theta^a}{N}\right)^N \xrightarrow{N \to \infty} e^{t_a \theta^a},
\]
which is the exponential of a matrix, defined by the power series
\[
    e^{t_a \theta^a} = \sum_{n=0}^{\infty} \frac{(t_a \theta^a)^n}{n!}.
\]
Thus, the exponential map provides a global parametrization of the Lie group \(G\) in terms of the elements of its Lie algebra \(\mathfrak{g}\) as follows:
\[
    \tilde{g}(\theta) = e^{t_a \theta^a} = 1 + \tilde{t}_a \theta^a + \frac{1}{2} \theta^a \theta^b \tilde{t}_{ab}  + \mathrm{O}(\theta^3),
\]
where
\[
    \tilde{t}_A = t_a, \quad \tilde{t}_{ab} = t_a t_b + t_b t_a.
\]

The same goes for the he group multiplication function \(m\) and the inversion function \(j\). We can define new functions \(\tilde{m}\), \(\tilde{j}\) as follows:
\[
    \begin{aligned}
        \tilde{j}(\theta) = - \theta^a,                                                    \\
        \tilde{g}(\tilde{j}(\theta)) = \tilde{g}(\theta)^{-1} = \dots = e^{-t_a \theta^a}, \\
    \end{aligned}
\]
while for the multiplication function we have the \textbf{Campbell-Baker-Hausdorff function} so that
\[
    \begin{aligned}
        [\dots ]
    \end{aligned}
\]

These new structure constants are not different from the original ones, since the exponential map is a reparametrization of the Lie group \(G\) in terms of the elements of its Lie algebra \(\mathfrak{g}\): \(\tilde{t}_a = t_a\) implies that
\[
    f^c_{\ ab} t_c = \left[ t_a,\,t_b \right] = \left[ \tilde{t}_a,\,\tilde{t}_b \right] = \tilde{f}^c_{\ ab} \tilde{t}_c = \tilde{f}^c_{\ ab} t_c,
\]
thus the structure constants of the Lie algebra \(\mathfrak{g}\) are the same in both parametrizations, as expected since the Lie algebra is an intrinsic property of the Lie group \(G\) and does not depend on the choice of parametrization:
\[
    f^c_{\ ab} = \tilde{f}^c_{\ ab}.
\]

Thus we have
\[
    \tilde{m}^a_{\ bc} = \frac{1}{2} \tilde{f}^a_{\ bc} = \frac{1}{2} f^a_{\ bc},
\]
and if we expand the CBH function in power series we find\footnote{There exists a closed form expression for the CBH function, but it is not important for our purposes: remind only that exact complicated  expressions for computing all the teerms exist, but we are only interested in the first few terms of the expansion.}
\[
    \tilde{m}^a(\theta, \theta') = \theta^a + \theta'^a + \frac{1}{2} f^a_{\ bc} \theta^b \theta'^c + \mathrm{O}(\theta^3, \theta'^3),
\]
which is like a polynomial expansion of the group multiplication function \(m\) in terms of the structure constants \(f^a_{\ bc}\) of the Lie algebra \(\mathfrak{g}\). The exponential map thus provides a global parametrization of the Lie group \(G\) in terms of the elements of its Lie algebra \(\mathfrak{g}\) and allows us to express the group multiplication function \(m\) in terms of the structure constants \(f^a_{\ bc}\) of the Lie algebra \(\mathfrak{g}\).

\begin{example}[Phase Group U(1)]
    Let us define the group of phase transformations U(1) as the group of complex numbers of unit modulus under multiplication:
    \[
        \mathrm{U}(1) = \{ z \in \mathbb{C} : |z| = 1 \} = \{ e^{i\theta} : \theta \in \mathbb{R} \}.
    \]
    It is an abelian Lie group of dimension 1 which represents the Gauge group of QED.
    The elements of U(1) can be represented by an angle \(\theta\) with an exponential parametrization as follows:
    \[
        \tilde{g}(\theta) = e^{i\theta^a t_a} = e^{i\theta},
    \]
    since the Lie algebra of U(1) is one-dimensional and spanned by the identity matrix \(t = 1\). The structure constants of the Lie algebra of U(1) are zero since the commutator of any two elements of the Lie algebra is zero:
    \[
        [t, t] = 0 \implies f_{ab}^{\ \ c} = 0.
    \]
\end{example}

\begin{example}[Special Unitary Group SU(2)]
    The special unitary group SU(2) is the group of \(2 \times 2\) unitary matrices with determinant 1:
    \[
        \mathrm{SU}(2) = \{ U \in \mathbb{C}^{2 \times 2} : U^\dagger U = \mathbb{I}_2, \det(U) = 1 \}.
    \]
    It is a non-abelian Lie group of dimension 3: \(\bs{\theta} \in \mathbb{R}^3\). The elements of SU(2) can be parametrized using the Pauli matrices \(\sigma_i\) as follows:
    \[
        \tilde{g}(\theta) = e^{i \theta^i \frac{\tau(\theta)}{2}},
    \]
    where
    \[
        \tau(\theta) = \begin{pmatrix}
            \theta^3              & \theta^1 - i \theta^2 \\
            \theta^1 + i \theta^2 & -\theta^3
        \end{pmatrix} = \theta^1 \begin{pmatrix}
            0 & 1 \\
            1 & 0
        \end{pmatrix} + \theta^2 \begin{pmatrix}
            0 & -i \\
            i & 0
        \end{pmatrix} + \theta^3 \begin{pmatrix}
            1 & 0  \\
            0 & -1
        \end{pmatrix} = \theta^i \sigma_i,
    \]
    Thus \(t_a = \frac{1}{2}\tau_a\) and the structure constants of the Lie algebra \(\mathfrak{su}(2)\) are given by the Levi-Civita symbol \(\epsilon_{ijk}\):
    \[
        [\sigma_i, \sigma_j] = 2i \epsilon_{ijk} \sigma_k \implies f_{ij}^{\ \ k} = 2 \epsilon_{ijk}.
    \]
\end{example}

\begin{example}[Special Orthogonal Group SO(3)]
    The special orthogonal group SO(3) is the group of \(3 \times 3\) real orthogonal matrices with determinant 1:
    \[
        \mathrm{SO}(3) = \{ R \in \mathbb{R}^{3 \times 3} : R^T R = \mathbb{I}_3, \det(R) = 1 \}.
    \]
    It is a non-abelian Lie group of dimension 3. The elements of SO(3) can be parametrized using the generators of rotations \(J_i\) as follows:
    \[
        \tilde{g}(\theta) = e^{i \theta^i J_i},
    \]
    where the generators \(J_i\) satisfy the commutation relations
    \[
        [J_i, J_j] = i \epsilon_{ijk} J_k.
    \]

    [-....-] to be finished
\end{example}